\interlude{One Perfect Moment}
It is nigh-impossible to discuss Puella Magi Madoka Magica for long without bringing up Goethe's Faust. Quoted directly in the form of graffiti of the German text appearing on a wall in Episode 2, Goethe's retelling of the ancient legend of a man who made a deal with the devil heavily informs the entire series.

A brief summary of Faust: Faust, an old man who is a wise sage but finds no joy in his life, makes a deal with the demon Mephistopheles to become young again and try living his life differently. Mephistopheles agrees to show Faust all the pleasures and joys of life he missed, but in return, if Faust ever experiences a moment of perfect happiness so great that he wishes to stop time and make it last forever, Faust will immediately die and go to Hell. The first part (published 1808, revised 1828) mostly follows Faust as he woos a young woman named Margarete (sometimes also known by the short form Gretchen). After he kills her brother, he leaves for a while to celebrate Walpurgisnacht, when German folklore says witches and demons have an orgy on Mt Brocken. He returns to find Gretchen is now mad and in prison, and she gave birth to his child but it was taken away. He tries to free her, but she is so delusional she cannot understand what is going on and he is forced to leave her behind as he flees the guards. \cite{ref94}

Part two (published 1832, the year of Goethe's death) is much stranger: Faust is now getting old again, a successful and wealthy man and a powerful sorcerer, and he has time-travel adventures, has an affair with Helen of Troy, and wins a war by bringing in an army of demons. At the end, he finally does something motivated solely by the good of another, instead of himself, and experiences a moment of perfect happiness. He dies, but because it was doing a good deed, he goes to judgment instead of immediately to Hell. Gretchen pleads with the Virgin Mary to let her guide him into Heaven, and Mary agrees. \cite{ref95}

References to Faust abound in the series. Beside the aforementioned graffiti, Faust quotes frequently appear as cryptograms inside the witches' barriers. More importantly, the story itself has many Faustian elements. Walpurgisnacht, for example, while it is referred to as an immensely powerful witch, appears to actually be an event involving many witches engaging in an orgy of destruction, just as in Faust. The witch's barriers are prisons created by overwriting reality with their own despair and madness, just like Gretchen experiences near the end of Faust's first part. A moment of perfect happiness leads directly to Hell for Faust, and this happens to multiple characters in the anime: Mami goes in moments from the blissful discovery that she has friends and allies to her brutal death; Kyoko's father is happy to have a congregation that listens to him, only to commit murder-suicide when he discovers how Kyoko made it happen; Sayaka experiences the happiness of knowing she has saved Kyousuke, only for that to turn out to be the beginning of her descent to despair and witchery.

Even moreso, the story of Madoka Magica is arguably a retelling of Faust. Kyubey is clearly Mephistopheles; he first appears as a cute animal, and is soon revealed as a frightening, powerful predator who offers wishes in exchange for souls. Just as Mephistopheles wants Faust to experience a moment of happiness and then descend forever into Hell, Kyubey is preying on the emotional highs and lows of the magical girls, and wants the energy released when they descend into despair and become witches.

Since Kyubey's primary target is Madoka, it might be tempting to see her as Faust, but that would be a mistake. The anime more readily compares her to Gretchen; her witch form is named Gretchen Kriemhilde, for example. Kyubey spends most of the anime trying and failing to get her to take the contract, before finally succeeding, just as Mephistopheles is frustrated in his first few attempts to corrupt Gretchen so that he can make her fall for Faust. Finally, Madoka's wish to guide magical girls away from being witches parallels Gretchen's wish to guide Faust into Heaven. Madoka also takes on a role as a savior and protector, similar to that taken in Faust by the Marian, divine principle of the Eternal Feminine, with which Gretchen is associated \cite{ref96}. 

If not Madoka, who is Faust? Homura is a fairly close match. Like Faust, she makes a bargain with the devil to turn back time and correct the mistakes she believes she has made. (More literally in Homura's case, but then again Faust eventually time-travels, too.) Her closeness to Madoka and desire to rescue her also reflect Faust's feelings for Gretchen, and her power to stop time may be a reference to the conditions of Faust's curse. Finally, like Faust she eventually learns that her attempt to turn back the clock has only made things worse.

However, Madoka Magica also subverts Faust. In the end, Homura's wish is not a mistake but key to breaking the cycle, and Madoka/Gretchen appeals to Kyubey/Mephistopheles, not Mary or God, to gain the power to guide others to Heaven. That is because Madoka is neither a character from Faust nor a Christian figure at all. Her true role is as a character from another mythology entirely.

Just as Madoka Magica presents a typically magical girl "false" veneer before pushing it aside to bring in the "true" form of the show, it also presents a Faustian, Christian surface reading that masks an underlying, essentially Buddhist story. Most immediately noticeably, one of the central tenets of Buddhism is that desire leads to suffering \cite{ref97}, and this is very much the case in Madoka Magica. All wishes lead ultimately to pain and despair; emotional highs are balanced by emotional lows. 

Another key Buddhist concept that the series addresses by name is karma. Karma is a very complex concept, and different sects view it very differently. Broadly, however, the Buddhist view can be very loosely summed up as cause and effect: action plants seeds which grow (maybe in this life, maybe in the next) into consequences. Good actions lead to good consequences and bad to bad, but either way, it has the effect of trapping you in the cycle of karma, because those consequences lead to further action which leads to more consequences \cite{ref98}. The magical girl happiness-despair cycle works in much the same way, dragging the magical girls steadily down to witch-hood. 

The weight of karma also binds people to a cycle of rebirth in Buddhist belief, forcing them to live over and over again, facing the burdens of the karma from past lives \cite{ref99}, much as Homura is trapped in the looping timelines illustrated in Episode 10. Enlightenment, the understanding of the true nature of the world, is the only way to escape this cycle of karma—and it is only on the last cycle that Madoka learns both of Homura's time travel (the cycle of rebirth) and precisely what the Incubators are doing (the nature of karma). Finally, Walpurgisnacht strongly resembles a lotus blossom (a symbol of enlightenment \cite{ref100}) while at the same time the gear motif reflects the ever-grinding wheel of dharma \cite{ref101}. 

As noted earlier, Madoka resembles a figure from Buddhist mythology, the bodhisattva Kanon (Japanese) or Kuanyin (Chinese). Kanon was a young girl who nearly attained nirvana, but stopped just before she reached it. She transcended space and time to reach out to others and help them to Enlightenment, before finally ascending to nirvana herself. This helps explain the Virgin Mary connection, as well, as the similarity between Kanon and Mary is quite noticeable and frequently commented upon \cite{ref102}. 

Thus it is that after saving everyone across time and space as a bodhisattva, Madoka then crosses the threshold to the next level. She becomes a force of nature, an incarnation of hope, dissolving her consciousness, and attaining nirvana.
Or that would be the plan, anyway. As we will see, she takes something of a detour along the way.
